Wasserstoffgas besteht aus $\mathrm{H_2}$-Molekülen. Mit der molaren Masse von Wasserstoff (FoTaBe S.~228) ist
\begin{align}
 M &= \MHyF = 2\cdot\MHO = \MHyg
\end{align}
Mit $T = \THF = \TH$ liefert die ideale Gasgleichung $pV = nRT$ die Stoffmenge
\begin{align}
 n &= \nHyF = \frac{\pH\cdot\VH}{\ncR\cdot\TH} = \nHy
\end{align}
Masse und Dichte sind damit
\begin{align}
 m &= \mHyF = \nHy\cdot\MHy = \mHyg \approx \result{\mHygP{0}{2}} \\
 \rho &= \rHyF = \frac{\mHy}{\VH} = \rHy \approx \result{\rHyP{0}{2}}
\end{align}
